\documentclass[12pt]{article}

\usepackage{float}
\usepackage{graphicx}			% Use this package to include images
\usepackage{amsmath}			% A library of many standard math expressions
\usepackage[margin=1in]{geometry}% Sets 1in margins. 
\usepackage{fancyhdr}			% Creates headers and footers
\usepackage{enumerate}          %These two package give custom labels to a list
\usepackage[shortlabels]{enumitem}
\usepackage{matlab-prettifier}
\usepackage{tikz}
\usepackage[utf8]{inputenc}
\usepackage{pgfplots} 
\usepackage{pgfgantt}
\usepackage{pdflscape}
\usepackage{array,multirow}
\usepackage{tabularx,booktabs}
\newcolumntype{Y}{>{\centering\arraybackslash}X}
\usetikzlibrary{shapes.geometric}
\usetikzlibrary{arrows.meta,arrows}
\pgfplotsset{compat=newest} 
\pgfplotsset{plot coordinates/math parser=false}
\newlength\fwidth
\newlength\fheight
\usepackage{float,graphicx}
\pagestyle{fancy}
\fancyhead[l]{}
\fancyhead[c]{Experiment 10}
\fancyhead[r]{}
\fancyfoot[c]{\thepage}
\setlength{\headheight}{20pt}
\begin{document}
\author{}
\date{}
\title{Thermal Expansion}
\maketitle
\section*{Aim}
To determine the coefficient of thermal expansion of various materials.
\section*{Analysis}
All materials expand when heated and contract when cooled.
For solids close to room temperature, the dependence is linear and can be described by the following equation.
\[\Delta L = L_0\alpha(T-T_0)\]
Rearranging terms we get
\[\frac{1}{\alpha} = L_0\cdot\frac{dT}{dL}\]
If $f$ is a product of $x_i$s then the uncertainty in $f$ is given by
\[\bigg(\frac{\sigma_f}{f}\bigg)^2=\sum_{i = 1}^{n}\bigg(\frac{\sigma_{x_i}}{x_i}\bigg)^2 \]
The error in $\alpha$ is given by
\[\sigma_\alpha = \sqrt{\bigg(\frac{\sigma_{L_0}}{L_0}\bigg)^2 +\bigg(\frac{\sigma^2_{m}}{m}\bigg)^2}\]
We calculate this for different materials.

\vspace*{7pt}

\noindent \textbf{Safety precautions :}
The induction stove must not go beyond 800 watts.
The pipes carring steam must not be bent.
Use gloves to touch hot equipment.
\subsection*{Brass}
\begin{figure}[H]
    \centering
    \setlength\fwidth{0.7\textwidth}
    \setlength\fheight{0.5\textwidth}
    % This file was created by matlab2tikz.
%
\definecolor{mycolor1}{rgb}{0.00000,0.44700,0.74100}%
\definecolor{mycolor2}{rgb}{0.85000,0.32500,0.09800}%
\definecolor{mycolor3}{rgb}{0.12941,0.12941,0.12941}%
%
\begin{tikzpicture}

\begin{axis}[%
width=0.951\fwidth,
height=\fheight,
at={(0\fwidth,0\fheight)},
scale only axis,
xmin=0,
xmax=0.8,
xlabel style={font=\color{mycolor3}},
xlabel={Change in length (in mm)},
ymin=20,
ymax=90,
ylabel style={font=\color{mycolor3}},
ylabel={Temperature (in C)},
axis background/.style={fill=white},
xmajorgrids,
ymajorgrids,
legend style={at={(0.03,0.97)}, anchor=north west, legend cell align=left, align=left}
]
\addplot[only marks, mark=o, mark options={}, mark size=1.5000pt, draw=mycolor1] table[row sep=crcr]{%
x	y\\
0.71	78\\
0.66	73\\
0.61	68\\
0.56	66\\
0.51	62\\
0.46	59\\
0.41	56\\
0.36	51\\
0.31	47\\
0.26	44\\
0.21	40\\
0.16	37\\
0.11	34\\
0.06	30\\
0.01	26\\
};
\addlegendentry{Data}

\addplot [color=mycolor2]
  table[row sep=crcr]{%
0	25.35143\\
0.00808080808080808	25.9361341616162\\
0.0161616161616162	26.5208383232323\\
0.0242424242424242	27.1055424848485\\
0.0323232323232323	27.6902466464646\\
0.0404040404040404	28.2749508080808\\
0.0484848484848485	28.859654969697\\
0.0565656565656566	29.4443591313131\\
0.0646464646464646	30.0290632929293\\
0.0727272727272727	30.6137674545455\\
0.0808080808080808	31.1984716161616\\
0.0888888888888889	31.7831757777778\\
0.096969696969697	32.3678799393939\\
0.105050505050505	32.9525841010101\\
0.113131313131313	33.5372882626263\\
0.121212121212121	34.1219924242424\\
0.129292929292929	34.7066965858586\\
0.137373737373737	35.2914007474748\\
0.145454545454545	35.8761049090909\\
0.153535353535354	36.4608090707071\\
0.161616161616162	37.0455132323232\\
0.16969696969697	37.6302173939394\\
0.177777777777778	38.2149215555556\\
0.185858585858586	38.7996257171717\\
0.193939393939394	39.3843298787879\\
0.202020202020202	39.969034040404\\
0.21010101010101	40.5537382020202\\
0.218181818181818	41.1384423636364\\
0.226262626262626	41.7231465252525\\
0.234343434343434	42.3078506868687\\
0.242424242424242	42.8925548484848\\
0.250505050505051	43.477259010101\\
0.258585858585859	44.0619631717172\\
0.266666666666667	44.6466673333333\\
0.274747474747475	45.2313714949495\\
0.282828282828283	45.8160756565657\\
0.290909090909091	46.4007798181818\\
0.298989898989899	46.985483979798\\
0.307070707070707	47.5701881414141\\
0.315151515151515	48.1548923030303\\
0.323232323232323	48.7395964646465\\
0.331313131313131	49.3243006262626\\
0.339393939393939	49.9090047878788\\
0.347474747474747	50.4937089494949\\
0.355555555555556	51.0784131111111\\
0.363636363636364	51.6631172727273\\
0.371717171717172	52.2478214343434\\
0.37979797979798	52.8325255959596\\
0.387878787878788	53.4172297575758\\
0.395959595959596	54.0019339191919\\
0.404040404040404	54.5866380808081\\
0.412121212121212	55.1713422424242\\
0.42020202020202	55.7560464040404\\
0.428282828282828	56.3407505656566\\
0.436363636363636	56.9254547272727\\
0.444444444444444	57.5101588888889\\
0.452525252525253	58.0948630505051\\
0.460606060606061	58.6795672121212\\
0.468686868686869	59.2642713737374\\
0.476767676767677	59.8489755353535\\
0.484848484848485	60.4336796969697\\
0.492929292929293	61.0183838585859\\
0.501010101010101	61.603088020202\\
0.509090909090909	62.1877921818182\\
0.517171717171717	62.7724963434343\\
0.525252525252525	63.3572005050505\\
0.533333333333333	63.9419046666667\\
0.541414141414141	64.5266088282828\\
0.54949494949495	65.111312989899\\
0.557575757575758	65.6960171515152\\
0.565656565656566	66.2807213131313\\
0.573737373737374	66.8654254747475\\
0.581818181818182	67.4501296363636\\
0.58989898989899	68.0348337979798\\
0.597979797979798	68.619537959596\\
0.606060606060606	69.2042421212121\\
0.614141414141414	69.7889462828283\\
0.622222222222222	70.3736504444445\\
0.63030303030303	70.9583546060606\\
0.638383838383838	71.5430587676768\\
0.646464646464647	72.1277629292929\\
0.654545454545455	72.7124670909091\\
0.662626262626263	73.2971712525253\\
0.670707070707071	73.8818754141414\\
0.678787878787879	74.4665795757576\\
0.686868686868687	75.0512837373737\\
0.694949494949495	75.6359878989899\\
0.703030303030303	76.2206920606061\\
0.711111111111111	76.8053962222222\\
0.719191919191919	77.3901003838384\\
0.727272727272727	77.9748045454545\\
0.735353535353535	78.5595087070707\\
0.743434343434344	79.1442128686869\\
0.751515151515152	79.728917030303\\
0.75959595959596	80.3136211919192\\
0.767676767676768	80.8983253535353\\
0.775757575757576	81.4830295151515\\
0.783838383838384	82.0677336767677\\
0.791919191919192	82.6524378383838\\
0.8	83.237142\\
};
\addlegendentry{Fit}

\end{axis}
\end{tikzpicture}%
    \caption{The slope is $71.36\pm0.85$}
\end{figure}
\[\alpha = \frac{1}{582\times 71.36} = 2.408\times 10^{-5}\]
The uncertainties are
\begin{align*}
    \sigma_{L_0} &= 1 \\
    \sigma_{slope} &= 0.85
\end{align*}
Thus we get
\[\sigma_\alpha = 1.203\% = 0.029\times10^{-5}\]
The coefficient of thermal expansion in brass is $2.408\pm 0.03 \times 10^{-5}$

\subsection*{Copper}
\begin{figure}[H]
    \centering
    \setlength\fwidth{0.7\textwidth}
    \setlength\fheight{0.5\textwidth}
    % This file was created by matlab2tikz.
%
\definecolor{mycolor1}{rgb}{0.00000,0.44700,0.74100}%
\definecolor{mycolor2}{rgb}{0.85000,0.32500,0.09800}%
\definecolor{mycolor3}{rgb}{0.12941,0.12941,0.12941}%
%
\begin{tikzpicture}

\begin{axis}[%
width=0.951\fwidth,
height=\fheight,
at={(0\fwidth,0\fheight)},
scale only axis,
xmin=0,
xmax=0.7,
xlabel style={font=\color{mycolor3}},
xlabel={Change in length (in mm)},
ymin=20,
ymax=100,
ylabel style={font=\color{mycolor3}},
ylabel={Temperature (in C)},
axis background/.style={fill=white},
xmajorgrids,
ymajorgrids,
legend style={at={(0.03,0.97)}, anchor=north west, legend cell align=left, align=left}
]
\addplot[only marks, mark=o, mark options={}, mark size=1.5000pt, draw=mycolor1] table[row sep=crcr]{%
x	y\\
0.68	90\\
0.63	86\\
0.58	81\\
0.53	77\\
0.48	73\\
0.43	69\\
0.38	64\\
0.33	59\\
0.28	55\\
0.23	51\\
0.18	45\\
0.13	41\\
0.08	35\\
0.03	30\\
};
\addlegendentry{Data}

\addplot [color=mycolor2]
  table[row sep=crcr]{%
0	28.65451\\
0.00707070707070707	29.3015962222222\\
0.0141414141414141	29.9486824444444\\
0.0212121212121212	30.5957686666667\\
0.0282828282828283	31.2428548888889\\
0.0353535353535354	31.8899411111111\\
0.0424242424242424	32.5370273333333\\
0.0494949494949495	33.1841135555556\\
0.0565656565656566	33.8311997777778\\
0.0636363636363636	34.478286\\
0.0707070707070707	35.1253722222222\\
0.0777777777777778	35.7724584444444\\
0.0848484848484848	36.4195446666667\\
0.0919191919191919	37.0666308888889\\
0.098989898989899	37.7137171111111\\
0.106060606060606	38.3608033333333\\
0.113131313131313	39.0078895555556\\
0.12020202020202	39.6549757777778\\
0.127272727272727	40.302062\\
0.134343434343434	40.9491482222222\\
0.141414141414141	41.5962344444444\\
0.148484848484848	42.2433206666667\\
0.155555555555556	42.8904068888889\\
0.162626262626263	43.5374931111111\\
0.16969696969697	44.1845793333333\\
0.176767676767677	44.8316655555556\\
0.183838383838384	45.4787517777778\\
0.190909090909091	46.125838\\
0.197979797979798	46.7729242222222\\
0.205050505050505	47.4200104444444\\
0.212121212121212	48.0670966666667\\
0.219191919191919	48.7141828888889\\
0.226262626262626	49.3612691111111\\
0.233333333333333	50.0083553333333\\
0.24040404040404	50.6554415555556\\
0.247474747474747	51.3025277777778\\
0.254545454545455	51.949614\\
0.261616161616162	52.5967002222222\\
0.268686868686869	53.2437864444444\\
0.275757575757576	53.8908726666667\\
0.282828282828283	54.5379588888889\\
0.28989898989899	55.1850451111111\\
0.296969696969697	55.8321313333333\\
0.304040404040404	56.4792175555556\\
0.311111111111111	57.1263037777778\\
0.318181818181818	57.77339\\
0.325252525252525	58.4204762222222\\
0.332323232323232	59.0675624444444\\
0.339393939393939	59.7146486666667\\
0.346464646464646	60.3617348888889\\
0.353535353535354	61.0088211111111\\
0.360606060606061	61.6559073333333\\
0.367676767676768	62.3029935555556\\
0.374747474747475	62.9500797777778\\
0.381818181818182	63.597166\\
0.388888888888889	64.2442522222222\\
0.395959595959596	64.8913384444444\\
0.403030303030303	65.5384246666667\\
0.41010101010101	66.1855108888889\\
0.417171717171717	66.8325971111111\\
0.424242424242424	67.4796833333333\\
0.431313131313131	68.1267695555556\\
0.438383838383838	68.7738557777778\\
0.445454545454545	69.420942\\
0.452525252525252	70.0680282222222\\
0.45959595959596	70.7151144444444\\
0.466666666666667	71.3622006666667\\
0.473737373737374	72.0092868888889\\
0.480808080808081	72.6563731111111\\
0.487878787878788	73.3034593333333\\
0.494949494949495	73.9505455555555\\
0.502020202020202	74.5976317777778\\
0.509090909090909	75.244718\\
0.516161616161616	75.8918042222222\\
0.523232323232323	76.5388904444444\\
0.53030303030303	77.1859766666667\\
0.537373737373737	77.8330628888889\\
0.544444444444444	78.4801491111111\\
0.551515151515151	79.1272353333333\\
0.558585858585859	79.7743215555555\\
0.565656565656566	80.4214077777778\\
0.572727272727273	81.068494\\
0.57979797979798	81.7155802222222\\
0.586868686868687	82.3626664444444\\
0.593939393939394	83.0097526666667\\
0.601010101010101	83.6568388888889\\
0.608080808080808	84.3039251111111\\
0.615151515151515	84.9510113333333\\
0.622222222222222	85.5980975555555\\
0.629292929292929	86.2451837777778\\
0.636363636363636	86.89227\\
0.643434343434343	87.5393562222222\\
0.65050505050505	88.1864424444444\\
0.657575757575757	88.8335286666667\\
0.664646464646465	89.4806148888889\\
0.671717171717172	90.1277011111111\\
0.678787878787879	90.7747873333333\\
0.685858585858586	91.4218735555555\\
0.692929292929293	92.0689597777778\\
0.7	92.716046\\
};
\addlegendentry{Fit}

\end{axis}
\end{tikzpicture}%
    \caption{The slope is $91.51\pm1.10$}
\end{figure}
\[\alpha = \frac{1}{582\times 91.51} = 1.878\times 10^{-5}\]
The uncertainties are
\begin{align*}
    \sigma_{L_0} &= 1 \\
    \sigma_{slope} &= 1.10
\end{align*}
Thus we get
\[\sigma_\alpha = 1.215\% = 0.023\times10^{-5}\]
The coefficient of thermal expansion in copper is $1.878\pm 0.023 \times 10^{-5}$

\subsection*{Aluminium}
\begin{figure}[H]
    \centering
    \setlength\fwidth{0.7\textwidth}
    \setlength\fheight{0.5\textwidth}
    % This file was created by matlab2tikz.
%
\definecolor{mycolor1}{rgb}{0.00000,0.44700,0.74100}%
\definecolor{mycolor2}{rgb}{0.85000,0.32500,0.09800}%
\definecolor{mycolor3}{rgb}{0.12941,0.12941,0.12941}%
%
\begin{tikzpicture}

\begin{axis}[%
width=0.951\fwidth,
height=\fheight,
at={(0\fwidth,0\fheight)},
scale only axis,
xmin=0,
xmax=0.8,
xlabel style={font=\color{mycolor3}},
xlabel={Change in length (in mm)},
ymin=20,
ymax=100,
ylabel style={font=\color{mycolor3}},
ylabel={Temperature (in C)},
axis background/.style={fill=white},
xmajorgrids,
ymajorgrids,
legend style={at={(0.03,0.97)}, anchor=north west, legend cell align=left, align=left}
]
\addplot[only marks, mark=o, mark options={}, mark size=1.5000pt, draw=mycolor1] table[row sep=crcr]{%
x	y\\
0.74	86\\
0.69	84\\
0.64	80\\
0.59	75\\
0.54	72\\
0.49	68\\
0.44	62\\
0.39	58\\
0.34	54\\
0.29	50\\
0.24	45\\
0.19	42\\
0.14	37\\
0.09	33\\
0.04	29\\
};
\addlegendentry{Data}

\addplot [color=mycolor2]
  table[row sep=crcr]{%
0	25.54548\\
0.00808080808080808	26.2248450909091\\
0.0161616161616162	26.9042101818182\\
0.0242424242424242	27.5835752727273\\
0.0323232323232323	28.2629403636364\\
0.0404040404040404	28.9423054545455\\
0.0484848484848485	29.6216705454545\\
0.0565656565656566	30.3010356363636\\
0.0646464646464646	30.9804007272727\\
0.0727272727272727	31.6597658181818\\
0.0808080808080808	32.3391309090909\\
0.0888888888888889	33.018496\\
0.096969696969697	33.6978610909091\\
0.105050505050505	34.3772261818182\\
0.113131313131313	35.0565912727273\\
0.121212121212121	35.7359563636364\\
0.129292929292929	36.4153214545455\\
0.137373737373737	37.0946865454546\\
0.145454545454545	37.7740516363636\\
0.153535353535354	38.4534167272727\\
0.161616161616162	39.1327818181818\\
0.16969696969697	39.8121469090909\\
0.177777777777778	40.491512\\
0.185858585858586	41.1708770909091\\
0.193939393939394	41.8502421818182\\
0.202020202020202	42.5296072727273\\
0.21010101010101	43.2089723636364\\
0.218181818181818	43.8883374545455\\
0.226262626262626	44.5677025454546\\
0.234343434343434	45.2470676363636\\
0.242424242424242	45.9264327272727\\
0.250505050505051	46.6057978181818\\
0.258585858585859	47.2851629090909\\
0.266666666666667	47.964528\\
0.274747474747475	48.6438930909091\\
0.282828282828283	49.3232581818182\\
0.290909090909091	50.0026232727273\\
0.298989898989899	50.6819883636364\\
0.307070707070707	51.3613534545455\\
0.315151515151515	52.0407185454546\\
0.323232323232323	52.7200836363636\\
0.331313131313131	53.3994487272727\\
0.339393939393939	54.0788138181818\\
0.347474747474747	54.7581789090909\\
0.355555555555556	55.437544\\
0.363636363636364	56.1169090909091\\
0.371717171717172	56.7962741818182\\
0.37979797979798	57.4756392727273\\
0.387878787878788	58.1550043636364\\
0.395959595959596	58.8343694545455\\
0.404040404040404	59.5137345454546\\
0.412121212121212	60.1930996363636\\
0.42020202020202	60.8724647272727\\
0.428282828282828	61.5518298181818\\
0.436363636363636	62.2311949090909\\
0.444444444444444	62.91056\\
0.452525252525253	63.5899250909091\\
0.460606060606061	64.2692901818182\\
0.468686868686869	64.9486552727273\\
0.476767676767677	65.6280203636364\\
0.484848484848485	66.3073854545455\\
0.492929292929293	66.9867505454546\\
0.501010101010101	67.6661156363636\\
0.509090909090909	68.3454807272727\\
0.517171717171717	69.0248458181818\\
0.525252525252525	69.7042109090909\\
0.533333333333333	70.383576\\
0.541414141414141	71.0629410909091\\
0.54949494949495	71.7423061818182\\
0.557575757575758	72.4216712727273\\
0.565656565656566	73.1010363636364\\
0.573737373737374	73.7804014545455\\
0.581818181818182	74.4597665454545\\
0.58989898989899	75.1391316363636\\
0.597979797979798	75.8184967272727\\
0.606060606060606	76.4978618181818\\
0.614141414141414	77.1772269090909\\
0.622222222222222	77.856592\\
0.63030303030303	78.5359570909091\\
0.638383838383838	79.2153221818182\\
0.646464646464647	79.8946872727273\\
0.654545454545455	80.5740523636364\\
0.662626262626263	81.2534174545455\\
0.670707070707071	81.9327825454546\\
0.678787878787879	82.6121476363636\\
0.686868686868687	83.2915127272727\\
0.694949494949495	83.9708778181818\\
0.703030303030303	84.6502429090909\\
0.711111111111111	85.329608\\
0.719191919191919	86.0089730909091\\
0.727272727272727	86.6883381818182\\
0.735353535353535	87.3677032727273\\
0.743434343434344	88.0470683636364\\
0.751515151515152	88.7264334545455\\
0.75959595959596	89.4057985454546\\
0.767676767676768	90.0851636363636\\
0.775757575757576	90.7645287272727\\
0.783838383838384	91.4438938181818\\
0.791919191919192	92.1232589090909\\
0.8	92.802624\\
};
\addlegendentry{Fit}

\end{axis}
\end{tikzpicture}%
    \caption{The slope is $84.07\pm0.92$}
\end{figure}
\[\alpha = \frac{1}{582\times 84.07} = 2.044\times 10^{-5}\]
The uncertainties are
\begin{align*}
    \sigma_{L_0} &= 1 \\
    \sigma_{slope} &= 0.92
\end{align*}
Thus we get
\[\sigma_\alpha = 1.110\% = 0.023\times10^{-5}\]
The coefficient of thermal expansion in aluminium is $2.044\pm 0.03 \times 10^{-5}$
\section*{Result}
The values of coefficient of thermal expansion for different materials is listed below.
\begin{align*}
    \text{Brass : }&2.408\pm 0.029 \times 10^{-5} \\
    \text{Copper : }&1.878\pm 0.023 \times 10^{-5} \\
    \text{Aluminium : }&1.110\pm 0.023 \times 10^{-5}
\end{align*}

\end{document}